% vim:ts=4:sw=4
%
% Copyright (c) 2008-2009 solvethis
% Copyright (c) 2010-2015 Casper Ti. Vector
% Public domain.
%
% 使用前请先仔细阅读 impthss 和 biblatex-caspervector 的文档，
% 特别是其中的 FAQ 部分和用红色强调的部分。
% 两者可在终端/命令提示符中用
%   texdoc impthss
%   texdoc biblatex-caspervector
% 调出。

% 采用了自定义的（包括大小写不同于原文件的）字体文件名，
% 并改动 ctex.cfg 等配置文件的用户请自行加入 nofonts 选项；
% 其它用户不用加入 nofonts 选项，加入之后反而会产生错误。
%
% 图书馆要求电子版论文的目录必须为黑色，
% 且某些教务要求打印版论文的文字部分为纯黑色而非灰度打印，
% 【因此最终打印和提交论文前，请将“colorlinks”改为“nocolorlinks”。】
\documentclass[UTF8, nocolorlinks, openany, spacing, pdftoc, pkuspace]{impthss}
% 设置行距为1.5。ctex说明文档中说linespread也可以作为选项传入，但测试结果显示不起效果。
% 因此在此处显示调用\linespread命令设置行距，测试显示有效果。
\linespread{1.5}
% 使用 biblatex 排版参考文献，并规定其格式（详见 biblatex-caspervector 的文档）。
% 这里按照英文文献在前，中文文献在后排序（“sorting = ecnty”）；
% 若需按照中文文献在前，英文文献在后排序，请设置“sorting = centy”；
% 若需按照引用顺序排序，请设置“sorting = none”。
\usepackage[backend = biber, style = caspervector, utf8, sorting = none]{biblatex}
% 提供近似于学校所要求的 Times New Roman / Arial 的字体。
\usepackage[defaultsups]{newtxtext}
\usepackage{newtxmath}
% 产生 originauth.tex 里的 \square。
\usepackage{latexsym}

% 按学校要求设定参考文献列表中的条目之内及之间的距离。
\setlength{\bibitemsep}{3bp}
% 对于 linespread 值的计算过程有兴趣的同学可以参考 impthss-extra.sty。
\renewcommand*{\bibfont}{\zihao{5}\linespread{1.27}\selectfont}

% 设定文档的基本信息。
% \ctexset{today=big}
\impthssinfo{
	cthesisname = {博士研究生学位论文}, ethesisname = {Doctor Thesis},
	ctitle = {中国科学院近代物理研究所博士论文模版},
	etitle = {Doctor Thesis \LaTeX{} Template for the Institute of Modern Physics of Chineses Academy of Sciences},
	cauthor = {周勇},
	eauthor = {Zhou Yong},
	studentid = {0123456789},
	date = {\today},
	school = {某某学院},
	cmajor = {粒子物理与原子核物理}, emajor = {Nuclear and Particle Physics},
	direction = {实验核物理},
	cmentor = {孙志宇教授}, ementor = {Prof.\ Sun Zhiyu},
	ckeywords = {其一，其二}, ekeywords = {KeywdOne, KeywdTwo}
}
% 载入参考文献数据库（注意不要省略“.bib”）。
\addbibresource{thesis.bib}

% extra packages
\usepackage[absolute, overlay]{textpos} %封面制作
\usepackage{import} % 使用\import,\subimport命令模块化文档结构，可以用相对路径载入文档
\usepackage{blindtext}
\usepackage{comment}
\usepackage{siunitx}

% 设置表格字体大小
\usepackage[font=normalsize]{caption}
\let\oldtabular\tabular 
\renewcommand{\tabular}{\small\oldtabular}

\begin{document}
	%%%%%%%%%%%%%%%%%%%%%%%%%
	%%% 封面,版权,声明,致谢 %%%
	\covermatter
	% 自动生成封面。
%	\maketitle
	% 从pdf载入封面
 	\begin{textblock*}{210mm}(0mm, 0mm)
\begin{figure*}[!hbt]
	\centering
	\includegraphics[width=210mm]{fig/ucas_cover_cn.eps}
\end{figure*}
\end{textblock*}

 	\null\newpage
 	\begin{textblock*}{210mm}(0mm, 0mm)
\begin{figure*}[!hbt]
	\centering
	\includegraphics[width=210mm]{fig/cover_en.eps}
\end{figure*}
\end{textblock*}


 	\null\newpage

	% 版权声明。
	% vim:ts=4:sw=4
%
% Copyright (c) 2008-2009 solvethis
% Copyright (c) 2010-2015 Casper Ti. Vector
% All rights reserved.
%
% Redistribution and use in source and binary forms, with or without
% modification, are permitted provided that the following conditions are
% met:
%
% * Redistributions of source code must retain the above copyright notice,
%   this list of conditions and the following disclaimer.
% * Redistributions in binary form must reproduce the above copyright
%   notice, this list of conditions and the following disclaimer in the
%   documentation and/or other materials provided with the distribution.
% * Neither the name of Peking University nor the names of its contributors
%   may be used to endorse or promote products derived from this software
%   without specific prior written permission.
%
% THIS SOFTWARE IS PROVIDED BY THE COPYRIGHT HOLDERS AND CONTRIBUTORS "AS
% IS" AND ANY EXPRESS OR IMPLIED WARRANTIES, INCLUDING, BUT NOT LIMITED TO,
% THE IMPLIED WARRANTIES OF MERCHANTABILITY AND FITNESS FOR A PARTICULAR
% PURPOSE ARE DISCLAIMED. IN NO EVENT SHALL THE COPYRIGHT HOLDER OR
% CONTRIBUTORS BE LIABLE FOR ANY DIRECT, INDIRECT, INCIDENTAL, SPECIAL,
% EXEMPLARY, OR CONSEQUENTIAL DAMAGES (INCLUDING, BUT NOT LIMITED TO,
% PROCUREMENT OF SUBSTITUTE GOODS OR SERVICES; LOSS OF USE, DATA, OR
% PROFITS; OR BUSINESS INTERRUPTION) HOWEVER CAUSED AND ON ANY THEORY OF
% LIABILITY, WHETHER IN CONTRACT, STRICT LIABILITY, OR TORT (INCLUDING
% NEGLIGENCE OR OTHERWISE) ARISING IN ANY WAY OUT OF THE USE OF THIS
% SOFTWARE, EVEN IF ADVISED OF THE POSSIBILITY OF SUCH DAMAGE.


\chapter*{版权声明}
% \markboth{版权声明}{} 
\thispagestyle{empty} %覆盖\chapter*{}中的调用的\thispagestyle{plain},因为此页不需要页眉页脚

任何收存和保管本论文各种版本的单位和个人，
未经本论文作者同意，不得将本论文转借他人，
亦不得随意复制、抄录、拍照或以任何方式传播。
否则一旦引起有碍作者著作权之问题，将可能承担法律责任。

% 若需排版二维码，请将二维码图片重命名为“barcode”，
% 转为合适的图片格式，并放在当前目录下，然后去掉下面 3 行的注释。
%\vfill\noindent
%\includegraphics[height = 5em]{barcode}


	% 原创性声明和使用授权说明。
	% vim:ts=4:sw=4
%
% Copyright (c) 2008-2009 solvethis
% Copyright (c) 2010-2015 Casper Ti. Vector
% All rights reserved.
%
% Redistribution and use in source and binary forms, with or without
% modification, are permitted provided that the following conditions are
% met:
%
% * Redistributions of source code must retain the above copyright notice,
%   this list of conditions and the following disclaimer.
% * Redistributions in binary form must reproduce the above copyright
%   notice, this list of conditions and the following disclaimer in the
%   documentation and/or other materials provided with the distribution.
% * Neither the name of Peking University nor the names of its contributors
%   may be used to endorse or promote products derived from this software
%   without specific prior written permission.
%
% THIS SOFTWARE IS PROVIDED BY THE COPYRIGHT HOLDERS AND CONTRIBUTORS "AS
% IS" AND ANY EXPRESS OR IMPLIED WARRANTIES, INCLUDING, BUT NOT LIMITED TO,
% THE IMPLIED WARRANTIES OF MERCHANTABILITY AND FITNESS FOR A PARTICULAR
% PURPOSE ARE DISCLAIMED. IN NO EVENT SHALL THE COPYRIGHT HOLDER OR
% CONTRIBUTORS BE LIABLE FOR ANY DIRECT, INDIRECT, INCIDENTAL, SPECIAL,
% EXEMPLARY, OR CONSEQUENTIAL DAMAGES (INCLUDING, BUT NOT LIMITED TO,
% PROCUREMENT OF SUBSTITUTE GOODS OR SERVICES; LOSS OF USE, DATA, OR
% PROFITS; OR BUSINESS INTERRUPTION) HOWEVER CAUSED AND ON ANY THEORY OF
% LIABILITY, WHETHER IN CONTRACT, STRICT LIABILITY, OR TORT (INCLUDING
% NEGLIGENCE OR OTHERWISE) ARISING IN ANY WAY OUT OF THE USE OF THIS
% SOFTWARE, EVEN IF ADVISED OF THE POSSIBILITY OF SUCH DAMAGE.

{
	\CTEXsetup[
		format+ = {\centering}, beforeskip = {40bp}, afterskip = {15bp}
	]{section}

	\chapter*{北京大学学位论文原创性声明和使用授权说明}
	\thispagestyle{empty} %覆盖\chapter*{}中的调用的\thispagestyle{plain},因为此页不需要页眉页脚
	\mbox{}\vspace*{-3em}
	\section*{原创性声明}

	本人郑重声明：
	所呈交的学位论文，是本人在导师的指导下，独立进行研究工作所取得的成果。
	除文中已经注明引用的内容外，
	本论文不含任何其他个人或集体已经发表或撰写过的作品或成果。
	对本文的研究做出重要贡献的个人和集体，均已在文中以明确方式标明。
	本声明的法律结果由本人承担。
	\vskip 1em
	\rightline{%
		论文作者签名：\hspace{5em}%
		日期：\hspace{2em}年\hspace{2em}月\hspace{2em}日%
	}

	\section*{%
		学位论文使用授权说明\\[-0.33em]
		\textmd{\zihao{5}（必须装订在提交学校图书馆的印刷本）}%
	}

	本人完全了解北京大学关于收集、保存、使用学位论文的规定，即：
	\begin{itemize}
		\item 按照学校要求提交学位论文的印刷本和电子版本；
		\item 学校有权保存学位论文的印刷本和电子版，
			并提供目录检索与阅览服务，在校园网上提供服务；
		\item 学校可以采用影印、缩印、数字化或其它复制手段保存论文；
		\item 因某种特殊原因需要延迟发布学位论文电子版，
			授权学校在 $\square$\nobreakspace{}一年 / %
			$\square$\nobreakspace{}两年 / %
			$\square$\nobreakspace{}三年以后在校园网上全文发布。
	\end{itemize}
	\centerline{（保密论文在解密后遵守此规定）}
	\vskip 1em
	\rightline{%
		论文作者签名：\hspace{5em}导师签名：\hspace{5em}%
		日期：\hspace{2em}年\hspace{2em}月\hspace{2em}日%
	}

}


	% 致谢。
	\subimport{chap/acknowledge/}{acknowledge}
	
	%%%%%%%%%%%%%%%%%%%%%
	%%% 中英文摘要,目录 %%%
	% 简单页眉页脚.此后内容都加入目录,并双面打印
	\frontmatter
	% 中英文摘要。
	\subimport{chap/abstract/}{abstract.tex}
	% 自动生成目录。	
	\tableofcontents

	%%%%%%%%%%%%%%%%%
	%%% 正文各章节 %%%
	\mainmatter
	% 各章节。
	\subimport{chap/introduction/}{introduction.tex}
	% 各章节。
	\subimport{chap/description/}{description.tex}
	\begin{comment}
	% 各章节
	\subimport{chap/dynamic_range/}{dynamic_range.tex}
	% 各章节。
	\subimport{chap/pmt_test/}{pmt_test.tex}
	% 各章节。
	\subimport{chap/construction/}{construction.tex}
	%
	\subimport{chap/cosmic_ray/}{cosmic_ray.tex}
	% 各章节。
	\subimport{chap/heavy_ion/}{heavy_ion.tex}
	% 结论
	\subimport{chap/conclusion/}{conclusion.tex}
	% special chapter没有章节编号，在目录中显示
% 	\subimport*{chap/special}{}
	\end{comment}
	% 正文中的附录部分。
	\appendix
	% 排版参考文献列表。
	\printbibliography[
		% 使“参考文献”出现在目录中；如果同时要使参考文献列表参与章节编号，
		% 可将“bibintoc”改为“bibnumbered”。
		heading = bibintoc,
		% 单独设定排序方案。此设定会局部覆盖之前的全局设置。
		% 注：只有同时使用 2.x 或之后版本的 biblatex 和相应兼容版本的 biber，
		% 才能对每个 \printbibliography 命令采用不同的排序方案，
		% 否则只能在载入 biblatex 宏包时就（全局）指定排序方案。
		% 在这样的情况下，请去掉所有的 sorting 选项，否则可能出错。
		% 此外，biblatex 3.0 中 \printbibliography 的 sorting 选项失效，
		% 详见 biblatex-caspervector 的文档。
		sorting = ecnty
	]
	% 各附录。
	\subimport*{chap/appendix/}{app1.tex}
	\subimport*{chap/appendix/}{app2.tex}
	
	
	% 以下为正文之后的部分，默认不进行章节编号。
	\backmatter
	
\end{document}

